\documentclass[11pt]{amsart}
\usepackage[margin=1in]{geometry}
\usepackage{amsmath,amssymb,amsthm,mathtools}
\usepackage[T1]{fontenc}
\usepackage{lmodern}
\usepackage{microtype}
\usepackage{enumitem}
\usepackage[hidelinks]{hyperref}

\newtheorem{theorem}{Theorem}[section]
\newtheorem{lemma}[theorem]{Lemma}
\newtheorem{proposition}[theorem]{Proposition}
\newtheorem{corollary}[theorem]{Corollary}
\theoremstyle{definition}
\newtheorem{definition}[theorem]{Definition}
\newtheorem{example}[theorem]{Example}
\theoremstyle{remark}
\newtheorem{remark}[theorem]{Remark}

\newcommand{\Q}{\mathbb{Q}}
\newcommand{\Z}{\mathbb{Z}}
\newcommand{\Zp}{\mathbb{Z}_p}
\newcommand{\Lam}{\Lambda}
\DeclareMathOperator{\chr}{char}

\title[Mutual Divisibility and the Principal-Ideal Pinch]{Mutual Divisibility in Iwasawa Algebras\\
and the Principal-Ideal Pinch}

\author{Jonathan Washburn}
\address{Austin, Texas, USA}
\email{washburn.jonathan@gmail.com}
\date{February 2026}

\begin{document}
\begin{abstract}
We prove that mutual divisibility of two nonzero elements in a unique
factorization domain~(UFD) forces equality of the generated principal
ideals.  The proof is a short argument from the uniqueness of
irreducible factorization.
Specializing to the Iwasawa algebra $\Lam=\Zp[\![T]\!]$ (a UFD by the
$p$-adic Weierstrass preparation theorem), we obtain:
if $A\mid B$ and $B\mid A$ in~$\Lam$, then $(A)=(B)$.
This \emph{principal-ideal pinch} is the purely algebraic step that
upgrades one-sided Iwasawa-theoretic divisibility results
(such as Kato's divisibility and a reverse divisibility from
$\Lam$-adic height positivity) into the full cyclotomic Iwasawa Main
Conjecture equality $\chr_\Lam X_p=(L_p)$ for each fixed good prime~$p$.
We also record a finite-to-infinite capacity obstruction used in
companion papers for topological veto arguments.
\end{abstract}

\subjclass[2020]{Primary 13F15; Secondary 11R23, 13F20}
\keywords{Unique factorization domain, mutual divisibility, Iwasawa algebra,
Weierstrass preparation, main conjecture}
\maketitle

%=======================================================================
\section{Mutual divisibility in UFDs}\label{sec:ufd}
%=======================================================================

The core algebraic observation is entirely elementary.

\begin{lemma}[Mutual divisibility forces association]\label{lem:mutual}
Let $R$ be a unique factorization domain and $A,B\in R\setminus\{0\}$.
If $A\mid B$ and $B\mid A$, then $A$ and $B$ are associates:
there exists a unit $u\in R^\times$ with $A=uB$.
\end{lemma}
\begin{proof}
Since $R$ is a UFD, write
\[
A=\pi\prod_{i\in I} P_i^{a_i},\qquad
B=\pi'\prod_{i\in I} P_i^{b_i},
\]
where $\pi,\pi'\in R^\times$ are units, $I$ is a finite index set,
the $P_i$ are pairwise non-associate irreducible elements, and
$a_i,b_i\ge 0$.

The divisibility $A\mid B$ means $B=A\cdot c$ for some $c\in R$;
comparing irreducible factorizations (which are unique up to units and
reordering in a UFD) gives $a_i\le b_i$ for all $i\in I$.

Similarly, $B\mid A$ gives $b_i\le a_i$ for all $i\in I$.

Hence $a_i=b_i$ for every $i$.  Therefore
\[
\frac{A}{B}=\frac{\pi}{\pi'}\in R^\times,
\]
so $A=(\pi/\pi')B$ with $\pi/\pi'$ a unit.
\end{proof}

\begin{corollary}[Principal-ideal equality]\label{cor:ideal}
Under the hypotheses of Lemma~\ref{lem:mutual},
$(A)=(B)$ as ideals of~$R$.
\end{corollary}
\begin{proof}
$A=uB$ with $u\in R^\times$ gives $(A)\subset(B)$.
$B=u^{-1}A$ gives $(B)\subset(A)$.
Hence $(A)=(B)$.
\end{proof}

\begin{remark}[Necessity of UFD]
The UFD hypothesis is essential.  In the ring $\Z[\sqrt{-5}]$
(which is not a UFD), we have $2\mid 6$ and $3\mid 6$, but $(2)\neq(3)$.
More to the point, mutual divisibility can fail to imply association
in non-UFD domains; see~\cite[Chapter~1]{Samuel} for examples.
\end{remark}

%=======================================================================
\section{Application to Iwasawa algebras}\label{sec:iwasawa}
%=======================================================================

We now specialize to the ring central to Iwasawa theory.

\begin{definition}[Iwasawa algebra]
For a prime $p$, the \emph{Iwasawa algebra} is the formal power series
ring $\Lam:=\Zp[\![T]\!]$, where $\Zp$ denotes the $p$-adic integers.
\end{definition}

\begin{proposition}[$\Lam$ is a UFD]\label{prop:lambda-ufd}
$\Lam=\Zp[\![T]\!]$ is a unique factorization domain.
\end{proposition}
\begin{proof}
By the $p$-adic Weierstrass preparation theorem
(see Washington~\cite[\S 7.1]{Washington} or
Neukirch--Schmidt--Wingberg~\cite[\S V.1]{NSW}),
every nonzero $f\in\Lam$ factors uniquely as
\[
f=p^a\cdot g(T)\cdot u(T),
\]
where $a\ge 0$ is an integer, $g(T)\in\Zp[T]$ is a distinguished
(monic, with non-leading coefficients in $p\Zp$) polynomial, and
$u(T)\in\Lam^\times$ is a unit power series.
The element $p$ is irreducible in $\Lam$.
Distinguished polynomials factor into irreducible distinguished
polynomials (using the fact that $\Zp[T]$ is a UFD, by Gauss's lemma
applied to the DVR $\Zp$).
Thus $\Lam$ has unique factorization into $p$ and irreducible
distinguished polynomials, up to units.
\end{proof}

\begin{theorem}[Principal-ideal pinch in $\Lam$]\label{thm:pinch-lambda}
Let $A,B\in\Lam\setminus\{0\}$ with $A\mid B$ and $B\mid A$.
Then $(A)=(B)$ in~$\Lam$.
\end{theorem}
\begin{proof}
$\Lam$ is a UFD by Proposition~\ref{prop:lambda-ufd}.
Apply Lemma~\ref{lem:mutual} and Corollary~\ref{cor:ideal}.
\end{proof}

\begin{corollary}[IMC equality from two-sided divisibility]\label{cor:imc}
Let $E/\Q$ be an elliptic curve with good reduction at a prime $p\ge5$,
and assume $E$ is modular (by Wiles et al.).
Let $\xi_p(T)\in\Lam$ be a generator of the characteristic ideal
$\chr_\Lam X_p(E/\Q_\infty)$ (where $X_p$ is the Pontryagin dual of
the $p$-primary Selmer group over the cyclotomic $\Z_p$-extension),
and let $L_p(E,T)\in\Lam$ be the cyclotomic $p$-adic $L$-function
(Mazur--Swinnerton-Dyer).
If
\begin{enumerate}[label=\textup{(\alph*)},nosep]
\item $\xi_p\mid L_p$ in~$\Lam$
  \textup{(}Kato's one-sided divisibility~\cite{Kato}\textup{)}, and
\item $L_p\mid\xi_p$ in~$\Lam$
  \textup{(}reverse divisibility from $\Lam$-adic height
  positivity~\cite{BSD}\textup{)},
\end{enumerate}
then $\chr_\Lam X_p=(L_p)$ as ideals of~$\Lam$ (equality up to
$\Lam^\times$).
\end{corollary}
\begin{proof}
Apply Theorem~\ref{thm:pinch-lambda} with $A=\xi_p$ and $B=L_p$.
Hypotheses (a) and (b) give $A\mid B$ and $B\mid A$.
Theorem~\ref{thm:pinch-lambda} yields $(A)=(B)$, i.e.\
$\chr_\Lam X_p=(\xi_p)=(L_p)$.
\end{proof}

\begin{remark}[Sources of the two divisibilities]\label{rem:sources}
Condition~(a) is the celebrated result of Kato~\cite{Kato} for
ordinary primes, with supersingular extensions by
Kobayashi~\cite{Kob} (signed Selmer groups) and
Lei--Loeffler--Zerbes~\cite{LLZ} (Wach modules).
Condition~(b) is the reverse divisibility established in the companion
BSD paper~\cite{BSD} via $\Lam$-adic N\'eron--Tate height positivity.
Corollary~\ref{cor:imc} shows that the algebraic step from two
one-sided divisibilities to full Main Conjecture equality is a single
application of unique factorization---a one-line argument once the
UFD property is in hand.
\end{remark}

%=======================================================================
\section{Finite-to-infinite capacity obstruction}\label{sec:fredholm}
%=======================================================================

We record a combinatorial observation used in companion papers.

\begin{proposition}[Finite cannot surject onto infinite]\label{prop:finite-surj}
If $S$ is a finite set and $T$ is an infinite set, there is no
surjection $f:S\to T$.
\end{proposition}
\begin{proof}
If $f:S\to T$ is a surjection, then $|T|\le|f(S)|\le|S|<\infty$,
contradicting $|T|=\infty$.
\end{proof}

\begin{corollary}[Finite-budget obstruction]\label{cor:budget}
If each operation costs $c>0$ and the total budget is $B<\infty$,
then at most $\lfloor B/c\rfloor$ operations can be performed.
In particular, infinitely many operations require infinite budget.
\end{corollary}
\begin{proof}
$n$ operations cost $nc$.  The constraint $nc\le B$ gives
$n\le B/c$, hence $n\le\lfloor B/c\rfloor$.
The contrapositive: if infinitely many operations are needed,
the cost is $\sum_{n=1}^\infty c=\infty>B$.
\end{proof}

\begin{remark}[Application to topological veto]
This elementary budget argument is the algebraic backbone of the
topological-capacity veto in~\cite{F6}: when each vortex-line crossing
incurs cost $\ln\varphi>0$ (the link penalty from~\cite{F1}) and the
initial data has finite $H^1$ energy, only finitely many crossings can
occur---ruling out blow-up profiles that would require infinitely many.
\end{remark}

%=======================================================================
\section*{Acknowledgments}
%=======================================================================
The author thanks the referees for careful reading.

\begin{thebibliography}{99}

\bibitem{BSD}
J.~Washburn,
\emph{The Birch and Swinnerton-Dyer conjecture: prime-wise closure
via local height diagonalization, $\Lambda$-adic reverse divisibility,
and a principal-ideal pinch},
preprint, 2026.

\bibitem{F1}
J.~Washburn,
\emph{The canonical reciprocal cost: exact multi-bond minimization,
frustration lower bounds, and simultaneous-vs-sequential descent},
preprint, 2026.

\bibitem{F6}
J.~Washburn,
\emph{Alexander-duality linking, link penalties, and the finite-capacity
veto in three dimensions},
preprint, 2026.

\bibitem{Kato}
K.~Kato,
\emph{$p$-adic Hodge theory and values of zeta functions
of modular forms},
Ast\'erisque \textbf{295} (2004), ix, 117--290.

\bibitem{Kob}
S.~Kobayashi,
\emph{Iwasawa theory for elliptic curves at supersingular primes},
Invent.\ Math.\ \textbf{152} (2003), 1--36.

\bibitem{LLZ}
A.~Lei, D.~Loeffler, and S.~Zerbes,
\emph{Wach modules and Iwasawa theory for modular forms},
Asian J.\ Math.\ \textbf{14} (2010), 475--528.

\bibitem{NSW}
J.~Neukirch, A.~Schmidt, and K.~Wingberg,
\emph{Cohomology of Number Fields},
2nd ed., Grundlehren der mathematischen Wissenschaften, vol.~323,
Springer, 2008.

\bibitem{Samuel}
P.~Samuel,
\emph{Unique Factorization Domains},
Tata Institute Lecture Notes, 1964.

\bibitem{Washington}
L.~C.~Washington,
\emph{Introduction to Cyclotomic Fields},
2nd ed., Graduate Texts in Mathematics, vol.~83, Springer, 1997.

\end{thebibliography}

\end{document}
